\documentclass[11pt]{article}
\usepackage{preamble}
\renewcommand{\answer}[1]{}

\begin{document}

\ladderheading{AP Statistics \textperiodcentered\ Unit 2 \textperiodcentered\ Topic 2.1 \textperiodcentered\ B: 2026-10-15 \textperiodcentered\ E: 2026-10-28}{Same Responses, Different Questions}

\focusbanner{TODAY'S PROBLEM}

Suppose a student council asks 120 ninth graders and 60 twelfth graders whether they prefer Field Day or a Talent Show for the spring celebration. The same responses are shown in two graphs.\par\smallskip \textbf{Andre:} \emph{Graph B hides how many people there are because both bars reach 100\%.}\par \textbf{Leila:} \emph{Graph A hides the relationship because the grade totals are different.}\par
\begin{center}
\begin{minipage}{0.48\linewidth}
\centering
\begin{tikzpicture}
\begin{axis}[
  width=0.95\linewidth,
  height=1.25in,
  ybar stacked,
  bar width=18pt,
  ymin=0,
  ymax=135,
  ytick={0,30,60,90,120},
  symbolic x coords={9th,12th},
  xtick=data,
  enlarge x limits=0.35,
  title={\textbf{Graph A}},
  xlabel={Grade},
  ylabel={Number of students},
  tick label style={font=\scriptsize},
  label style={font=\scriptsize},
  title style={font=\small},
  ymajorgrids,
  nodes near coords,
  every node near coord/.append style={font=\scriptsize},
  legend style={at={(0.5,-0.30)},anchor=north,font=\scriptsize,draw=none}
]
\addplot[fill=calloutblue,draw=dayblue] coordinates {(9th,72) (12th,18)};
\addplot[fill=calloutgreen,draw=dayblue] coordinates {(9th,48) (12th,42)};
\legend{Field Day,Talent Show}
\end{axis}
\end{tikzpicture}
\end{minipage}\hfill
\begin{minipage}{0.48\linewidth}
\centering
\begin{tikzpicture}
\begin{axis}[
  width=0.95\linewidth,
  height=1.25in,
  ybar stacked,
  bar width=18pt,
  ymin=0,
  ymax=108,
  ytick={0,20,40,60,80,100},
  symbolic x coords={9th,12th},
  xtick=data,
  enlarge x limits=0.35,
  title={\textbf{Graph B}},
  xlabel={Grade},
  ylabel={Percent within grade},
  tick label style={font=\scriptsize},
  label style={font=\scriptsize},
  title style={font=\small},
  ymajorgrids,
  nodes near coords,
  every node near coord/.append style={font=\scriptsize},
  legend style={at={(0.5,-0.30)},anchor=north,font=\scriptsize,draw=none}
]
\addplot[fill=calloutblue,draw=dayblue] coordinates {(9th,60) (12th,30)};
\addplot[fill=calloutgreen,draw=dayblue] coordinates {(9th,40) (12th,70)};
\legend{Field Day,Talent Show}
\end{axis}
\end{tikzpicture}
\end{minipage}
\end{center}
\begin{firsttakebox}{FIRST TAKE --- before the video (5 min)}
For the question \emph{Is celebration preference related to grade level?}, which graph would you use first? Name one visible feature.
\workspace{0.9in}
\end{firsttakebox}

\begin{callout}{\IconBook\ MATCH THE DISPLAY TO THE QUESTION (keep on the page)}{calloutgreen}
A \textbf{count bar graph} preserves how many observations are in each category and group. A
\textbf{segmented bar graph} stacks each group's conditional relative frequencies to $100\%$.
Equal total heights remove the effect of unequal group sizes, so corresponding segments compare proportions.
Different conditional distributions indicate an association. A segmented graph needs group totals before it
can recover counts; a count graph needs division before it directly shows within-group percentages.
\end{callout}

\newpage
\textbf{Finish after the video:}\par\smallskip

\dokbadge{1}~\textbf{(a)} From Graph A, state each grade's total surveyed and the count in each preference category.\par
\workspace{0.7in}

\dokbadge{2}~\textbf{(b)} Compute the conditional distribution of preference within each grade. Describe the largest contrast between the grades.\par
\workspace{1.4in}

\dokbadge{3}~\textbf{(c)} For deciding whether preference is related to grade level, choose the more useful graph and explain how equal-height bars settle the comparison. State what Graph A is useful for. Then give one question that Graph A cannot answer \emph{directly from bar lengths alone} and one question that Graph B cannot answer from the display.\par
\workspace{2.0in}

\begin{sentenceframebox}\raggedright\textbf{Frame:} Graph \blank[0.4in] better answers the relationship question because equal-height bars let me compare \blank[2.2in].\end{sentenceframebox}

\begin{sentenceframebox}\raggedright\textbf{Frame:} Graph A preserves \blank[1.2in], while Graph B cannot tell me \blank[2.0in].\end{sentenceframebox}

\begin{summaryexitbox}{EXIT --- turn this in}
One thing the video changed about my first take: \hrulefill\par
\workspace{0.4in}
\end{summaryexitbox}

\end{document}
